\emph{(i)} With $\nInf = \nu + e_{-}$ and $\nO = \tfrac12(\nu - e_{-})$, $\nu = \tfrac12\nInf + \nO$, so $Z = e_{0}\nu = \tfrac12 e_{0}\nInf + e_{0}\nO = P_{e_{0}} + e_{0}\nO$; and $Z^{2} = -e_{0}^{2}\nu^{2} = -1$. $Z$ commutes with the six bivectors $e_{i}e_{j}$ and $e_{i}e_{-}$ ($1 \le i < j \le 3$), which with $Z$ span the maximal compact subalgebra $\so(2)\oplus\so(4)$, and not with the boost $e_{0}e_{1}$; that $Z$ spans the centre of that subalgebra is Proposition \claimref{synthesis.positive-energy-selects-rule-2-sign}(i). For $a \perp \nO, \nInf$, $P_{a}^{2} = \tfrac14 a\nInf a\nInf = -\tfrac14 a^{2}\nInf^{2} = 0$, so $\exp(tP_{a}) = 1 + tP_{a}$ (the null case of \claimref{std-algebra.bivector-exponential}), and the check confirms that it maps $X(x)$ to $X(x + ta)$; more generally $P_{a}P_{b} = -\tfrac14 ab\,\nInf^{2} = 0$, so the translations commute. In the chart of Lemma \claimref{mass.chart-form-spin-isomorphism} the translations are the $x\wedge\nInf$ and the dilation is $\nInf\wedge\nO$; here $D_{\rho}$, with bivector $e_{-}\nu$ of square $1$, sends $\nInf \mapsto \rho^{2}\nInf$ and $\nO \mapsto \rho^{-2}\nO$ and fixes $\R^{1,3}$, so $X(x) \mapsto \rho^{-2}X(\rho^{2}x)$ and $D_{\rho}Z\tilde D_{\rho} = \rho^{2}P_{e_{0}} + \rho^{-2}e_{0}\nO$. The boost $B_{\rho}$ sends $e_{0} \mapsto \cosh\lambda\,e_{0} + \sinh\lambda\,e_{1}$ with $e^{\lambda} = \rho^{2}$ and fixes $\nInf$, so $B_{\rho}P_{e_{0}}\tilde B_{\rho} = P_{\cosh\lambda e_{0} + \sinh\lambda e_{1}}$, and $2\rho^{-2}(\cosh\lambda\,e_{0} + \sinh\lambda\,e_{1}) = (1 + \rho^{-4})e_{0} + (1 - \rho^{-4})e_{1} = \ell + \rho^{-4}(e_{0} - e_{1})$. For a rational unit $\hat n \in \mathrm{span}\{e_{1}, e_{2}, e_{3}\}$, the versor $V = (e_{1} + \hat n)e_{1}$ (two reflections) has $V\tilde V$ a nonzero scalar and $V e_{1}\tilde V/(V\tilde V) = \hat n$, fixing $e_{0}$ and $\nInf$; rational directions are dense in the sphere, and the general unit direction is reached by a rotation in the same way. A future timelike $a = (a_{0}, \vec a)$, $a_{0} > |\vec a| > 0$, is $\alpha(e_{0} + \hat n) + \beta(e_{0} - \hat n)$ with $\hat n = \vec a/|\vec a|$, $\alpha = \tfrac12(a_{0} + |\vec a|) > 0$, $\beta = \tfrac12(a_{0} - |\vec a|) > 0$ (and $a = a_{0}e_{0}$ is the half-sum of two opposite null vectors). The check verifies all of this exactly, with the coefficients tabulated: $\rho = 2, \tfrac32, 10$ for the dilation, the remainders $16^{-k}$ at $\rho = 2^{k}$, $(\cosh\lambda, \sinh\lambda) = (\tfrac{17}{8}, \tfrac{15}{8})$ and $(\tfrac{41}{9}, \tfrac{40}{9})$ at $\rho = 2, 3$, three directions, and the decompositions $(5,3,0,0) = 4\ell + (e_{0} - e_{1})$ and $(9, 2, -3, 6) = 8(e_{0} + \hat n) + (e_{0} - \hat n)$, $\hat n = \tfrac17(2, -3, 6)$.

\emph{(ii)} Let $\mathcal{D}$ be the Garding space of $\pi$, on which every $-i\,d\pi(Y)$ is essentially self-adjoint and which $\pi(g)$ preserves, with $d\pi(\mathrm{Ad}_{g}Y) = \pi(g)\,d\pi(Y)\,\pi(g)^{-1}$ (standard, Segal; cited from memory). For fixed $v \in \mathcal{D}$, $Y \mapsto \ip{v}{-i\,d\pi(Y)v}$ is linear on the finite-dimensional Lie algebra, hence continuous, so $C_{\pi} = \bigcap_{v}\{Y : \ip{v}{-i\,d\pi(Y)v} \ge 0\}$ is an intersection of closed half-spaces: a closed convex cone; it is invariant because $\ip{v}{-i\,d\pi(\mathrm{Ad}_{g}Y)v} = \ip{\pi(g)^{-1}v}{-i\,d\pi(Y)\pi(g)^{-1}v}$; and positivity on the core $\mathcal{D}$ gives positivity of the self-adjoint closure. Suppose $Z \in C_{\pi}$. By invariance and the cone property, $\rho^{-2}\mathrm{Ad}_{D_{\rho}}Z = P_{e_{0}} + \rho^{-4}e_{0}\nO \in C_{\pi}$ for every $\rho > 0$; letting $\rho \to \infty$, closedness gives $P_{e_{0}} \in C_{\pi}$. Then $2\rho^{-2}\mathrm{Ad}_{B_{\rho}}P_{e_{0}} = P_{\ell} + \rho^{-4}P_{e_{0} - e_{1}} \in C_{\pi}$, and $P_{\ell} \in C_{\pi}$ in the limit. Rotations give $P_{e_{0} + \hat n} \in C_{\pi}$ for every unit $\hat n$, and convexity with linearity of $a \mapsto P_{a}$ gives $P_{a} \in C_{\pi}$ for every $a$ in the closed forward cone $\overline{V}_{+}$. The operators $-i\,d\pi(P_{a})$ commute (the $P_{a}$ do and the translation subgroup is abelian), so they have a joint spectral measure on momenta $p$ with $-i\,d\pi(P_{a}) = \ip{a}{p}$ spectrally; $\ip{a}{p} \ge 0$ for all $a \in \overline{V}_{+}$ means $p \in \overline{V}_{+}$, the forward cone being self-dual. In particular every future null translation has a nonnegative generator. With $-Z$ in place of $Z$ the same steps give $-P_{a} \in C_{\pi}$ for $a \in \overline{V}_{+}$, the backward cone; which of the two holds is the choice of time orientation. A spacelike $a$ can be rotated to $|a|e_{1}$, and $\mathrm{Ad}_{e_{1}e_{2}}P_{e_{1}} = -P_{e_{1}}$; if $P_{a} \in C_{\pi}$ then also $-P_{a} \in C_{\pi}$, so $-i\,d\pi(P_{a})$ is both $\ge 0$ and $\le 0$, hence $0$. The check verifies the conjugation of $P_{e_{1}}$ to its negative, that the same rotation fixes $P_{e_{0}}$, and that the $-Z$ orbit tends to $-P_{e_{0}}$.

\emph{(iii)} Let $\pi$ act on the GNS space of the state and let its GNS vector $\Omega$ satisfy $\pi(\exp P_{a})\Omega = \Omega$ for all $a$ (for a second-quantised field, the Fock vacuum of the filled-sea representation, on which the second-quantised generators are $d\Gamma$ of the positive-energy one-particle ones and are nonnegative when those are; the hypothesis is not applied to a one-particle space such as $L^{2}(\R, du)$ of a generator, where the null momentum has both signs). A local Rindler horizon of Minkowski space through $p$ with generator direction $\ell$ is a null plane $\{p + u\ell + y\}$, and the translation along its generators is $\exp(tP_{\ell})$, the same for every generator. By (ii) its generator is nonnegative, and $\Omega$ is invariant: $\Omega$ is a ground state for the null translations of every such plane. That is the hypothesis, positivity and invariance, under which the result filed alongside, synthesis.positive-null-energy-gives-r-premises, shows that the modes of the plane on either side of the edge are jointly KMS, for the local boost future-directed on that side, at rapidity period $2\pi$; the local boost acts on each generator as the dilation about its edge point, and the wedge reflection, the geometric part of the modular conjugation in Fact \claimref{angle.imported-theorems-deliver-their-hypotheses}, is $\mathrm{Ad}(e_{0}e_{1})$, which sends $X(x_{0}, x_{1}, x_{2}, x_{3})$ to $-X(-x_{0}, -x_{1}, x_{2}, x_{3})$, the same null ray (verified exactly). The general version of the same conclusion, on the full wedge algebra, is the Bisognano-Wichmann theorem of that Fact, whose spectrum-condition hypothesis is what (ii) supplies. Minkowski space here is flat spacetime with vanishing gravitational field. The Schwarzschild field's field strength is pure nonzero Weyl, so no gauge makes it a pure dilation of the background and the background conformal transformations are not conformal symmetries of it; (ii) says nothing about the affine translations along its horizon's generators. Those along a cosmological horizon of the de Sitter field $\hb = \sigma$, a pure dilation, are background conformal transformations: for $\Lambda > 0$, $I = \nInf + \tfrac{\Lambda}{6}\nO$ and a future null $n$ with $\ip{n}{I} = 0$, in the physical affine parameter whose physical tangent is $n$ the uniform affine translation along the generators is generated by $P_{n} + \tfrac{\Lambda}{12}\,n\nO$, and the orbit identities of (i), applied to $e_{0}\nO$ with the limits taken as in (ii), place $n\nO$, hence that generator, in $C_{\pi}$, which gives the positivity there for conformally invariant fields; this is proved in the result filed alongside, synthesis.de-sitter-horizon-translation-is-conformal-null-rotation. The scope of (iii) as stated here is the local Rindler horizons of Minkowski space.
